% Minimal preamble for P2 cards taken from sources/comparability-partitions-2408.13523v2/source0.tex
% (arXiv:2408.13523v2, "Partitioning perfect graphs into comparability graphs").
% Every line below is copied mechanically from that file; nothing was invented, renamed or redefined.
%   line 2    : \documentclass[12pt]{article}
%   line 3    : \usepackage{amssymb,latexsym}
%   line 4    : \usepackage{amsmath}
%   line 5    : \usepackage[dvipdfm]{graphicx}  -- the dvipdfm driver option is dropped, because the
%               assembler forces the xelatex engine and a forced driver option contradicts it.
%   line 16   : \usepackage{amsfonts}
%   line 17   : \usepackage{amsthm}
%   lines 20-30: the author's theorem declarations
%   lines 32-41: the author's own macro definitions
% Omitted source preamble lines and why they are not needed for these body spans:
%   line 1 (provenance comment), lines 6-11 (tikz and its libraries: no drawing in the retained spans),
%   lines 7 and 9 (verbatim), lines 12-13 (page geometry), line 14 (rawfonts),
%   line 15 (psfrag), line 18 (hyperref with the pagebackref option: the retained spans print no
%   citation anchors and the assembly adds its own hyperref machinery), lines 43-61 (the author's
%   backref neutralisation and his margin-note macros, which depend on the omitted hyperref).
% The retained spans of the body use only \em, \bf, \cal, amsthm environments, amsmath displays and
% the macros defined in lines 32-41 below.
\documentclass[12pt]{article}
\usepackage{amssymb,latexsym}
\usepackage{amsmath}
\usepackage{graphicx}
\usepackage{amsfonts}
\usepackage{amsthm}

\newtheorem{theorem}{Theorem}%[section]
\newtheorem{problem}[theorem]{Problem}%[section]
\newtheorem{observation}[theorem]{Observation}
\newtheorem{lemma}[theorem]{Lemma}
\newtheorem{corollary}[theorem]{Corollary}
\newtheorem{definition}{Definition}
\newtheorem{example}[theorem]{Example}
\newtheorem{conjecture}{Conjecture}%[section]
\newtheorem{proposition}[theorem]{Proposition}
\newtheorem{claim}[theorem]{Claim}
\newtheorem{remark}[theorem]{Remark}

\def\qed{\ifhmode\unskip\nobreak\fi\quad\ifmmode\Box\else$\Box$\fi}
\def\pont{\hspace{-6pt}{\bf.\ }}
\def\qed{\hfill $\Box$}

\def\R{{\mbox{\boldmath$\mathcal R$}}}
% the author's second definition of \R, left commented out exactly as in the source
\def\F{{\mathcal F}}
\def\P{{\mathcal P}}

\def\newpage{\vfill\eject}
\title{Almost all perfect graphs have their edge set partitioned into two comparability graphs}
\author{András Gyárfás, Márton Marits, Géza Tóth}
\date{arxiv:2408.13523v2; CC BY 4.0}
\begin{document}
\maketitle
\noindent\textit{Editorial note.} The counted unit is the article's own Theorem twocover at source lines 190-191 together with the author's proof, whose header at source line 246 names the theorem as Proof of Theorem twocover; the proof is not adjacent to the statement, because the article places the discussion of almost all, the definition of GSP graphs, the theorem of Promel and Steger and the theorem on GSP graphs in between, so the pairing rests on the author's own named header rather than on a later-first-proof rule.\par\medskip
\noindent\textit{Editorial note.} The retained context at source lines 117-120 and 128-129 is the article's own description of comparability graphs and its definitions of p(G) and c(G); the retained context at source lines 193-197, 204-206, 212-213 and 214-216 is its own explanation of the term almost all, its definition of GSP graphs and its statement of the theorem of Promel and Steger with its limiting meaning.\par\medskip
\noindent\textit{Editorial note.} The retained context at source lines 225-226 is the article's own theorem that every GSP graph has partition number at most two, which the proof uses as the inclusion of GSP(n) in P2PE(n). It is retained as a prerequisite; its own proof is not part of this material.\par\medskip
\noindent\textit{Editorial note.} The retained context at source lines 212-213 carries the author's citation to the paper of Promel and Steger, and the bibliography fragment attached to this task carries that single entry copied verbatim from the author's own bibliography at source line 750; no entry was rewritten, added or reordered.\par\medskip
\noindent\textit{Editorial note.} The source preamble cannot be retained as it stands: it loads rawfonts, psfrag, verbatim, tikz with several libraries, hyperref with the pagebackref option and a dvi driver option for graphicx, none of which the retained body spans use, and it redefines backref commands that only exist because of that hyperref option. The retained preamble file therefore carries, line for line, source line 2, source line 3, source line 4, source line 16, source line 17, the theorem declarations of source lines 20-30 and the author macro definitions of source lines 32-41; the only textual change anywhere is the removal of the dvipdfm driver option from the graphicx line, because the assembly forces the xelatex engine. No new macro, package or declaration was introduced.\par\medskip
\noindent\textit{Editorial note.} The article's Theorem promelsteger and Theorem gensplit are cited here only through their statements in the retained context, their proofs are not retained, and the article's other results, namely Theorem bipline, Theorem unit, Theorem co-triangulated, Theorem loglog, Theorem ub, Lemma rec, the corollaries genupbound, shift, doblesh and p', and Proposition H9, are not counted here.\par\medskip
\noindent\textit{Allowed standard prerequisites (editorial).}\begin{itemize}
\item The partition number p(G) and the cover number c(G). p(G) is the minimum m such that the edge set of G is the union of m edge-disjoint comparability subgraphs of G, and c(G) is the same number without the edge-disjointness requirement, so that c(G) is at most p(G); retained at source lines 128-129, with comparability graphs described at source lines 117-120.. The counted theorem asserts c(G) <= p(G) <= 2 for almost all perfect graphs, so both parameters and their relation must be fixed.
\item The density classes PE(n), P2PE(n), BE(n) and GSP(n) and the meaning of almost all. PE(n) is the set of labelled perfect graphs on n vertices and P2PE(n) the subset that can be partitioned into at most two comparability graphs, so that the statement of the counted theorem is the limit of the ratios |P2PE(n)|/|PE(n)| equal to one; BE(n) and GSP(n) are the corresponding sets of labelled Berge graphs and labelled GSP graphs; retained at source lines 193-197 and 214-216.. Supplies the exact quantitative form of the two almost-all statements that the proof multiplies, and the identification PE(n)=BE(n) by the strong perfect graph theorem that the proof invokes.
\item Generalized split graphs and the theorem of Promel and Steger. A graph is a GSP graph when either it or its complement splits into a complete graph and a union of vertex-disjoint complete graphs with no edges between them, and almost all Berge graphs are GSP, meaning the limit of |GSP(n)|/|BE(n)| equals one; retained at source lines 204-206 and 212-216.. Provides the class whose partition number the proof needs to bound and the density statement that transfers that bound to all Berge graphs.
\item Every GSP graph has partition number at most two. For any GSP graph G the edge set can be partitioned into two comparability graphs of G; retained at source lines 225-226.. Converts membership in GSP(n) into membership in P2PE(n), which is the inclusion the proof uses for the first limit.
\end{itemize}\medskip
Here we focus on a well known family of perfect graphs, the {\em comparability graphs},
which can be obtained from a partially ordered set as follows.
Vertices correspond to the elements of the partially ordered set,
and two vertices are connected by an undirected edge if and only if the corresponding elements

Let $p(G)$ denote the minimum number $m$ such that $E(G)=\cup_{i=1}^m E(G_i)$ where the $G_i$-s are edge-disjoint comparability subgraphs of $G$.
Dropping  the condition that the $G_i$-s are edge-disjoint, we define $c(G)$ as the smallest number of comparability subgraphs of $G$ needed to {\em cover} $E(G)$.

In Theorem \ref{twocover} we use the term ``almost all'' in the following sense.
Let $PE(n)$ denote the set of labelled perfect graphs on $n$ vertices and let $P2PE(n)$ denote the
set of labelled perfect graphs on $n$ vertices that can be partitioned into (at most) two comparability graphs.
Then  Theorem \ref{twocover} states that
$$\lim_{n\rightarrow \infty} {|P2PE(n)|\over |PE(n)|}=1.$$

\begin{definition} A graph $G$ is a GSP if either $G$ or its complement can be written in the following form:
  the vertex set is partitioned into $V_1,V_2$ where $V_1$ is a complete graph and $V_2$ is the union of vertex disjoint complete graphs with no edges between them.
\end{definition}

\begin{theorem}\label{promelsteger}\cite{PS} Almost all Berge graphs are GSP.
\end{theorem}

As in Theorem \ref{twocover}, ``almost all'' in Theorem \ref{promelsteger} means
$$\lim_{n\rightarrow \infty} { |GSP(n)|\over|BE(n)|}=1,$$
where $BE(n),GSP(n)$ denote the set of all labelled Berge graphs (respectively GSP graphs) on $n$ vertices.

\begin{theorem} \label{gensplit} For any GSP graph $G$, $p(G)\le 2$.
\end{theorem}

\begin{theorem} \label{twocover} We have $c(G)\le p(G)\le 2$ for almost all perfect graphs.
\end{theorem}

{\bf Proof of Theorem \ref{twocover}.} Using Theorem  \ref{gensplit} and the strong perfect graph theorem, we have
$GSP(n)\subset P2PE(n)$ and $BE(n)=PE(n)$, therefore from Theorem \ref{promelsteger},

$$1=\lim_{n\rightarrow \infty} {|GSP(n)|\over |BE(n)|}\le \lim_{n\rightarrow \infty}{|P2PE(n)|\over |PE(n)|}\le 1,$$
implying that the second limit is also equal to $1$, finishing the proof. \qed

\begin{thebibliography}{99}
\bibitem{PS} H. J. Pr\"omel, A. Steger, Almost all Berge graphs are perfect, {\em Combinatorics, Probability and Computing} {\bf 1.} (1992), 53--79.
\end{thebibliography}
\end{document}
